\documentclass[aspectratio=169,professionalfonts]{beamer}

% =================================================
% Idioma y codificación
% =================================================
\usepackage[spanish,es-nodecimaldot]{babel}
\usepackage[utf8]{inputenc}
\usepackage[T1]{fontenc}
\usepackage{lmodern}
\usepackage{microtype}

% =================================================
% Notas para presentador
% =================================================
\setbeameroption{show notes on second screen=left}
% \setbeameroption{show notes on second screen=right}

% =================================================
% Paquetes
% =================================================
\usepackage{hyperref}
\usepackage{graphicx}
\usepackage{booktabs}
\usepackage{array}
\usepackage{xcolor}
\usepackage{listings}
\usepackage{tikz}
\usetikzlibrary{positioning}
\usepackage{pdfcomment}
\usepackage{enumitem}
\usepackage{caption}
\usepackage{ragged2e}

% =================================================
% Colores
% =================================================
\definecolor{PVNavy}{HTML}{0F172A}
\definecolor{PVBlue}{HTML}{1D4ED8}
\definecolor{PVSky}{HTML}{38BDF8}
\definecolor{PVTeal}{HTML}{0F766E}
\definecolor{PVGold}{HTML}{D4A017}
\definecolor{PVRed}{HTML}{DC2626}
\definecolor{PVGreen}{HTML}{059669}
\definecolor{PVSlate}{HTML}{334155}
\definecolor{PVLight}{HTML}{F8FAFC}
\definecolor{PVGray}{HTML}{E2E8F0}
\definecolor{PVDarkText}{HTML}{0B1324}

% =================================================
% Tema visual elegante
% =================================================
\setbeamercolor{normal text}{fg=PVDarkText,bg=white}
\setbeamercolor{structure}{fg=PVBlue}
\setbeamercolor{title}{fg=white}
\setbeamercolor{subtitle}{fg=white}
\setbeamercolor{frametitle}{fg=white,bg=PVNavy}
\setbeamercolor{block title}{fg=white,bg=PVBlue}
\setbeamercolor{block body}{fg=PVDarkText,bg=PVLight}
\setbeamercolor{alertblock title}{fg=white,bg=PVRed}
\setbeamercolor{alertblock body}{fg=PVDarkText,bg=white}
\setbeamercolor{exampleblock title}{fg=white,bg=PVTeal}
\setbeamercolor{exampleblock body}{fg=PVDarkText,bg=white}

\setbeamertemplate{navigation symbols}{}
\setbeamertemplate{headline}{}
\setbeamertemplate{footline}{%
  \leavevmode
  \hbox{%
    \begin{beamercolorbox}[wd=.78\paperwidth,ht=3.1ex,dp=1.2ex,leftskip=1em]{author in head/foot}%
      \color{PVSlate}\small\insertshorttitle
    \end{beamercolorbox}%
    \begin{beamercolorbox}[wd=.22\paperwidth,ht=3.1ex,dp=1.2ex,rightskip=1em plus1fil]{date in head/foot}%
      \color{PVSlate}\small\insertframenumber/\inserttotalframenumber
    \end{beamercolorbox}%
  }
}

\setbeamertemplate{itemize item}{\color{PVBlue}$\blacktriangleright$}
\setbeamertemplate{itemize subitem}{\color{PVTeal}$\bullet$}

\setlength{\leftmargini}{1.2em}
\setlist[itemize]{itemsep=0.35em}
\setlist[enumerate]{itemsep=0.35em}

% =================================================
% Listings
% =================================================
\lstset{
  basicstyle=\ttfamily\small,
  breaklines=true,
  columns=fullflexible,
  frame=single,
  framerule=0.5pt,
  rulecolor=\color{black!20},
  backgroundcolor=\color{black!2},
  showstringspaces=false,
  keywordstyle=\color{PVBlue}\bfseries,
  commentstyle=\color{PVTeal},
  stringstyle=\color{PVRed}
}

% =================================================
% Macros
% =================================================
\newcommand{\bluebf}[1]{\textcolor{PVBlue}{\textbf{#1}}}
\newcommand{\soft}[1]{\textcolor{PVSlate}{#1}}

\newcommand{\pnote}[1]{%
  \pdfcomment[
    icon=Note,
    open=false,
    subject={Presenter Note},
    color={1 1 0},
    opacity=0,
    author={Presenter Viewer}
  ]{#1}%
}

\newcommand{\sectioncard}[2]{
\begin{frame}[plain]
\begin{tikzpicture}[remember picture,overlay]
  \fill[PVNavy] (current page.south west) rectangle (current page.north east);
  \fill[PVBlue] (current page.south west) rectangle ([xshift=0.22\paperwidth]current page.north west);
  \fill[PVSky,opacity=.15] ([xshift=0.55\paperwidth,yshift=0.15\paperheight]current page.south west) circle (2.4cm);
  \fill[PVSky,opacity=.10] ([xshift=0.8\paperwidth,yshift=0.75\paperheight]current page.south west) circle (3.0cm);
\end{tikzpicture}

\vspace*{1.8cm}
{\color{white}\fontsize{28}{32}\selectfont\bfseries #1\par}
\vspace{0.35cm}
{\color{white!85}\Large #2\par}
\vfill
{\color{white!65}\small Presenter Viewer Demo}
\end{frame}
}

\newcommand{\kbd}[1]{%
\tikz[baseline=(key.base)]\node[
  draw=PVGray,
  fill=white,
  rounded corners=3pt,
  inner xsep=5pt,
  inner ysep=2pt,
  text=PVNavy,
  font=\ttfamily\footnotesize
] (key) {#1};%
}

\newcommand{\shortcutline}[2]{
\begin{block}{#1}
\centering #2
\end{block}
}

\newcommand{\imgplaceholder}[2]{
\begin{center}
\IfFileExists{images/#1}{%
  \includegraphics[width=0.92\linewidth]{images/#1}%
}{%
  \fbox{\parbox[c][0.26\paperheight][c]{0.86\linewidth}{%
    \centering
    \Large Captura pendiente\\[0.4em]
    \texttt{images/#1}\\[0.6em]
    \normalsize #2
  }}%
}
\end{center}
}

\newcommand{\dashboardcard}[3]{
\begin{beamercolorbox}[wd=\linewidth,sep=0.9em,rounded=true,shadow=true]{block body}
{\color{PVBlue}\large\bfseries #1}\\[0.35em]
{\normalsize #2}\\[0.45em]
{\color{PVSlate}\small #3}
\end{beamercolorbox}
}

\newcommand{\codebox}[1]{%
\fbox{%
\parbox{0.94\linewidth}{%
\ttfamily\small #1
}%
}%
}

% =================================================
% Portada personalizada
% =================================================
\setbeamertemplate{title page}{
\begin{tikzpicture}[remember picture,overlay]
  \fill[PVNavy] (current page.south west) rectangle (current page.north east);
  \fill[PVBlue] (current page.south west) rectangle ([xshift=0.25\paperwidth]current page.north west);
  \fill[PVSky,opacity=.14] ([xshift=0.82\paperwidth,yshift=0.72\paperheight]current page.south west) circle (2.7cm);
  \fill[PVSky,opacity=.10] ([xshift=0.72\paperwidth,yshift=0.22\paperheight]current page.south west) circle (3.3cm);
  \fill[white,opacity=.06] ([xshift=0.34\paperwidth,yshift=0.18\paperheight]current page.south west)
       rectangle ([xshift=0.92\paperwidth,yshift=0.82\paperheight]current page.south west);
\end{tikzpicture}

\vspace*{0.72cm}
\begin{columns}[T,totalwidth=\textwidth]
\column{0.56\textwidth}
{\color{white!85}\large Demo elegante y tutorial}\par
\vspace{0.35cm}
{\color{white}\fontsize{26}{30}\selectfont\bfseries \inserttitle\par}
\vspace{0.35cm}
{\color{white!80}\large \insertsubtitle\par}
\vspace{0.9cm}
{\color{white!95}\normalsize \insertauthor\par}
\vspace{0.18cm}
{\color{white!72}\fontsize{8.5}{10}\selectfont \insertdate\par}

\column{0.39\textwidth}
\vspace{0.45cm}
\begin{beamercolorbox}[wd=\linewidth,sep=1.1em,rounded=true,shadow=true]{block body}
{\large\bfseries Contenido}\par
\vspace{0.5em}
\begin{itemize}
  \item Instalación y ejecución
  \item Interfaz del visor
  \item Teclas y herramientas
  \item \texttt{\textbackslash note\{\}} y \texttt{\textbackslash pnote\{\}}
  \item Flujo recomendado de trabajo
\end{itemize}
\end{beamercolorbox}
\end{columns}
\vfill
}

% =================================================
% Datos
% =================================================
\title{Presenter Viewer}
\subtitle{Manual de uso y preparación de presentaciones}
\author{Alan Díaz Manríquez (calix35@gmail.com)}
\date{\today}

\begin{document}

\begin{frame}[plain]
  \titlepage
  \note{Esta demo debe servir como archivo oficial de ejemplo y como PDF de prueba del proyecto.}
\end{frame}

\begin{frame}{Contenido}
\tableofcontents
\note{Recorrer rápidamente la estructura del manual antes de entrar a los detalles.}
\end{frame}

\section{Visión general}

\begin{frame}{¿Qué es Presenter Viewer?}
\begin{columns}[T,totalwidth=\textwidth]
\column{0.58\textwidth}
\begin{block}{Definición}
Presenter Viewer es un \bluebf{visor de presentaciones para el expositor} diseñado para trabajar cómodamente con PDFs académicos o profesionales, mostrando en una misma interfaz la información que el ponente necesita durante la exposición.
\end{block}

\begin{itemize}
  \item Visualiza la \bluebf{filmina actual}.
  \item Permite consultar la \bluebf{siguiente filmina}.
  \item Soporta \bluebf{notes} exportadas desde Beamer.
  \item Integra \bluebf{pnotes} como apoyo breve y táctico.
  \item Incluye herramientas interactivas para presentar en vivo.
\end{itemize}

\column{0.38\textwidth}
\begin{exampleblock}{Idea clave}
No es solo un lector de PDF: busca ofrecer \bluebf{control total de la presentación} sin que el público vea la lógica interna del expositor.
\end{exampleblock}
\end{columns}

\note{
Presentar Presenter Viewer como una herramienta de exposición, no solo como visor.
La diferencia importante es que separa lo que ve el público de lo que necesita el presentador.
}
\end{frame}

\begin{frame}{¿Para qué sirve?}

\begin{columns}[T,totalwidth=\textwidth]
\column{0.48\textwidth}
\begin{small}
\dashboardcard{Clases}{Permite exponer con notes, cronómetro, navegación visual y apoyo privado.}{Útil para cursos, laboratorios, demostraciones y sesiones híbridas.}
\end{small}
\vspace{0.3cm}
\begin{small}
\dashboardcard{Conferencias}{Facilita mantener ritmo, foco visual y separación limpia entre presentador y proyector.}{Especialmente valioso cuando hay dos pantallas o proyector externo.}
\end{small}
\column{0.48\textwidth}
\begin{small}
  \dashboardcard{Defensas}{Ayuda a resaltar detalles, controlar tiempos y conservar recordatorios discretos.}{Muy útil en tesis, coloquios, seminarios y presentaciones técnicas.}
\end{small}

\vspace{0.3cm}
\begin{small}
\dashboardcard{Ensayos}{Sirve como entorno realista de práctica antes de la exposición formal.}{Permite validar flujo, herramientas, layout y tiempo de intervención.}
\end{small}
\end{columns}

\note{
Aquí conviene enfatizar que la herramienta sirve tanto en escenario real como en preparación previa.
La idea es que el usuario ensaye en el mismo entorno donde luego va a presentar.
}
\end{frame}

\begin{frame}{Flujo de trabajo general}
\begin{block}{Pipeline recomendado}
\centering
\kbd{LaTeX} \hspace{0.35cm} $\rightarrow$ \hspace{0.35cm}
\kbd{PDF} \hspace{0.35cm} $\rightarrow$ \hspace{0.35cm}
\kbd{Presenter Viewer} \hspace{0.35cm} $\rightarrow$ \hspace{0.35cm}
\kbd{Configurar paneles} \hspace{0.35cm} $\rightarrow$ \hspace{0.35cm}
\kbd{Presentar}
\end{block}

\vspace{0.25cm}

\begin{columns}[T,totalwidth=\textwidth]
\column{0.58\textwidth}
\begin{itemize}
  \item Crear la presentación en Beamer.
  \item Exportar el PDF final.
  \item Abrirlo en Presenter Viewer.
  \item Ajustar paneles y herramientas.
  \item Conectar el proyector o segunda pantalla.
  \item Exponer con control del flujo.
\end{itemize}

\column{0.38\textwidth}
\begin{alertblock}{Objetivo}
Separar claramente la \bluebf{vista del presentador} de la \bluebf{salida del público}.
\end{alertblock}
\end{columns}

\note{
Esta diapositiva debe dejar claro el pipeline completo.
Presenter Viewer aparece como puente entre el PDF ya compilado y la exposición real.
}
\end{frame}

\section{Instalación y ejecución}
\sectioncard{1. Instalación y ejecución}{Cómo iniciar Presenter Viewer correctamente}

% =================================================
% Requisitos
% =================================================
\begin{frame}{Requisitos}
\begin{columns}[T,totalwidth=\textwidth]
\column{0.58\textwidth}
\begin{block}{Requisitos generales}
\begin{itemize}
  \item Python 3.x
  \item Dependencias del proyecto
  \item Entorno gráfico compatible con PyQt
  \item Un archivo PDF listo para abrir
\end{itemize}
\end{block}

\begin{exampleblock}{Recomendación}
Trabaja dentro de un \bluebf{entorno virtual} para mantener las dependencias aisladas y evitar conflictos con otras instalaciones de Python.
\end{exampleblock}

\column{0.38\textwidth}
\dashboardcard{Consejo}{Prueba primero con un PDF simple y después con uno que incluya notes.}{Así validas tanto el flujo básico como el modo de presentador.}
\end{columns}

\note{
Aclarar que la aplicación es gráfica, por lo que necesita un entorno con soporte de ventana.
}
\end{frame}

% =================================================
% Clonar repo
% =================================================
\begin{frame}[fragile]{Clonar el repositorio}
\begin{block}{Paso 1}
\begin{lstlisting}
git clone https://github.com/calix35/PresenterViewer.git
cd presenter_viewer
\end{lstlisting}
\end{block}

\begin{exampleblock}{Buena práctica}
Mantén el proyecto en una carpeta dedicada para que también sea sencillo ubicar \texttt{samples}, \texttt{images} y los archivos PDF de prueba.
\end{exampleblock}

\note{
Este es el punto de arranque para cualquier usuario que quiera instalar el sistema desde el repositorio.
}
\end{frame}

% =================================================
% Entorno virtual
% =================================================
\begin{frame}[fragile]{Crear y activar un entorno virtual}
\begin{columns}[T,totalwidth=\textwidth]
\column{0.58\textwidth}
\begin{block}{Crear el entorno}
\begin{lstlisting}
python -m venv .venv
\end{lstlisting}
\end{block}

\begin{block}{Activar en Windows}
\begin{lstlisting}
.venv\Scripts\activate
\end{lstlisting}
\end{block}

\begin{block}{Activar en Linux/macOS}
\begin{lstlisting}
source .venv/bin/activate
\end{lstlisting}
\end{block}

\column{0.38\textwidth}
\begin{alertblock}{Importante}
Conviene activar el entorno antes de instalar dependencias o ejecutar la aplicación.
\end{alertblock}
\end{columns}

\note{
Aquí puedes explicar que esto evita instalar paquetes globalmente.
}
\end{frame}

% =================================================
% requirements y dependencias
% =================================================
\begin{frame}[fragile]{Instalar dependencias desde \texttt{requirements.txt}}
\begin{block}{Instalación básica}
\begin{lstlisting}
python -m pip install -r requirements.txt
\end{lstlisting}
\end{block}

\begin{block}{Dependencia clave}
Entre las dependencias debe encontrarse \texttt{PyQt} o la variante de Qt que use el proyecto, ya que la interfaz gráfica depende de ella.
\end{block}

\begin{exampleblock}{¿Por qué usar \texttt{requirements.txt}?}
Porque permite reproducir el entorno exacto necesario para ejecutar Presenter Viewer en otra computadora.
\end{exampleblock}

\note{
Enfatizar que aquí vive la base del entorno gráfico: sin Qt, la aplicación no abre correctamente.
}
\end{frame}

% =================================================
% editable
% =================================================
\begin{frame}[fragile]{Instalación del proyecto (modo editable recomendado)}
\begin{columns}[T,totalwidth=\textwidth]
\column{0.58\textwidth}
\begin{block}{Instalación}
\begin{lstlisting}
python -m pip install -e .
\end{lstlisting}
\end{block}

\begin{itemize}
  \item Registra el comando \texttt{presenter-viewer}.
  \item Facilita ejecutar la aplicación como programa.
  \item Permite modificar el código sin reinstalar todo cada vez.
\end{itemize}

\column{0.38\textwidth}
\begin{alertblock}{Recomendado para desarrollo}
Este modo es especialmente útil cuando todavía se están haciendo pruebas, ajustes o correcciones.
\end{alertblock}
\end{columns}

\note{
Aquí puedes mencionar que el modo editable fue importante para simplificar la ejecución durante el desarrollo.
}
\end{frame}

% =================================================
% ejecución
% =================================================
\begin{frame}[fragile]{Ejecución de la aplicación}
\begin{block}{Comando principal}
\begin{lstlisting}
presenter-viewer
\end{lstlisting}
\end{block}

\begin{block}{Abrir un PDF directamente}
\begin{lstlisting}
presenter-viewer archivo.pdf
\end{lstlisting}
\end{block}

\begin{block}{Alternativa}
\begin{lstlisting}
python -m presenter_viewer.main
\end{lstlisting}
\end{block}

\pnote{Conviene iniciar primero con un PDF de prueba para verificar la detección de notes y paneles.}

\note{
Comentar que el archivo puede darse con ruta relativa o absoluta.
}
\end{frame}

% =================================================
% comportamiento inicial
% =================================================
\begin{frame}{Comportamiento al iniciar}
\begin{block}{Escenarios habituales}
\begin{itemize}
  \item Si se pasa un PDF al iniciar, la aplicación lo abre automáticamente.
  \item Si no se pasa PDF, puede intentar cargar un archivo de ejemplo.
  \item Si no existe un ejemplo, la ventana queda lista para abrir un archivo manualmente.
\end{itemize}
\end{block}

\begin{alertblock}{Recomendación}
Antes de una exposición real, abre explícitamente el PDF final y confirma que cargó correctamente.
\end{alertblock}

\note{
Esto ayuda a que el usuario entienda qué esperar incluso si arranca la aplicación vacía.
}
\end{frame}

% =================================================
% sin PDF
% =================================================
\begin{frame}{Aplicación iniciada sin PDF}
\begin{columns}[T,totalwidth=\textwidth]
\column{0.56\textwidth}
\imgplaceholder{startup-empty.png}{Ventana principal recién abierta, sin archivo cargado.}

\column{0.40\textwidth}
\begin{block}{Qué debería verse}
\begin{itemize}
  \item ventana principal activa;
  \item paneles vacíos o en espera;
  \item controles disponibles;
  \item opción clara para abrir un archivo.
\end{itemize}
\end{block}

\begin{exampleblock}{Interpretación}
Este estado indica que la aplicación está lista para trabajar, aunque todavía no se haya cargado ningún PDF.
\end{exampleblock}
\end{columns}

\note{
Esta captura ayuda mucho al usuario nuevo, porque le muestra que una ventana vacía no significa error.
}
\end{frame}

% =================================================
% con PDF
% =================================================
\begin{frame}{Aplicación iniciada con PDF cargado}
\begin{columns}[T,totalwidth=\textwidth]
\column{0.56\textwidth}
\imgplaceholder{startup-loaded.png}{Ventana principal con un PDF ya cargado desde línea de comandos o desde la interfaz.}

\column{0.40\textwidth}
\begin{block}{Qué debería verse}
\begin{itemize}
  \item filmina actual;
  \item paneles auxiliares activos;
  \item notas si el PDF las incluye;
  \item interfaz lista para presentar.
\end{itemize}
\end{block}

\begin{alertblock}{Verificación rápida}
Conviene revisar de inmediato si el archivo fue detectado como PDF normal o como PDF con notes.
\end{alertblock}
\end{columns}

\note{
Esta diapositiva conecta la instalación con el uso real: ya no solo corre la app, sino que queda lista para exponer.
}
\end{frame}

% =================================================
% abrir manualmente
% =================================================
\begin{frame}{Abrir un PDF manualmente}
\begin{columns}[T,totalwidth=\textwidth]
\column{0.58\textwidth}
\begin{block}{Métodos de apertura}
\begin{itemize}
  \item desde el botón o menú de abrir archivo;
  \item con el atajo de teclado;
  \item arrastrando el archivo PDF a la ventana.
\end{itemize}
\end{block}

\begin{block}{Atajo sugerido}
\begin{itemize}
  \item \kbd{Ctrl+O} en Windows y Linux
  \item \kbd{Cmd+O} en macOS
\end{itemize}
\end{block}

\column{0.38\textwidth}
\begin{exampleblock}{Ventaja}
Esto permite abrir rápidamente un PDF incluso aunque la aplicación haya iniciado sin archivo.
\end{exampleblock}
\end{columns}

\note{
Aquí conviene mencionar explícitamente el atajo según plataforma.
}
\end{frame}

% =================================================
% drag and drop
% =================================================
\begin{frame}{Soporte de drag \& drop}
\begin{columns}[T,totalwidth=\textwidth]
\column{0.58\textwidth}
\begin{block}{Flujo}
El usuario puede arrastrar un archivo PDF desde el explorador o gestor de archivos y soltarlo directamente sobre la ventana de Presenter Viewer.
\end{block}

\begin{itemize}
  \item agiliza la apertura;
  \item evita navegar por cuadros de diálogo;
  \item hace más natural el uso en escritorio.
\end{itemize}

\column{0.38\textwidth}
\begin{alertblock}{Importante}
El archivo debe ser un PDF válido y conviene verificar de inmediato si se detectó como PDF normal o como PDF con notes.
\end{alertblock}
\end{columns}

\note{
Este slide documenta una función muy cómoda y moderna que vale la pena destacar.
}
\end{frame}

\section{Uso básico del visor}
\sectioncard{2. Uso básico del visor}{Cómo trabajar con la interfaz durante una exposición}

% =================================================
% Interfaz principal
% =================================================
\begin{frame}{Interfaz principal}
\begin{columns}[T,totalwidth=\textwidth]
\column{0.56\textwidth}
\imgplaceholder{presenter-view.png}{Ventana principal del presentador con PDF cargado, paneles y barra inferior.}

\column{0.40\textwidth}
\begin{block}{Elementos principales}
\begin{itemize}
  \item panel principal (filmina actual);
  \item paneles auxiliares;
  \item barra inferior de estado;
  \item herramientas activas;
  \item cronómetro y controles.
\end{itemize}
\end{block}

\begin{exampleblock}{Idea clave}
El presentador ve múltiples elementos, pero el público solo ve la filmina.
\end{exampleblock}
\end{columns}

\note{
Esta es la vista más importante del sistema. Aquí el usuario debe entender la diferencia entre lo que él ve y lo que ve el público.
}
\end{frame}

% =================================================
% Barra de herramientas
% =================================================
\begin{frame}{Barra de herramientas}
\begin{columns}[T,totalwidth=\textwidth]
\column{0.56\textwidth}
\imgplaceholder{toolbar.png}{Barra superior con botones de control y herramientas del visor.}

\column{0.40\textwidth}
\begin{block}{Elementos comunes}
\begin{itemize}
  \item abrir PDF;
  \item navegación básica;
  \item modos de interacción;
  \item acceso rápido a funciones clave.
\end{itemize}
\end{block}

\begin{exampleblock}{Idea clave}
La barra concentra las acciones más usadas durante la presentación.
\end{exampleblock}
\end{columns}

\note{
Aquí puedes explicar cada botón rápidamente. No entrar a detalle, eso va en la sección de herramientas.
}
\end{frame}

% =================================================
% Distribución de paneles
% =================================================
\begin{frame}{Distribución de paneles}
\begin{columns}[T,totalwidth=\textwidth]
\column{0.58\textwidth}
\begin{block}{Distribución recomendada}
\begin{itemize}
  \item panel grande: filmina actual;
  \item panel superior derecho: siguiente;
  \item panel medio derecho: notas actuales;
  \item panel inferior derecho: notas siguientes.
\end{itemize}
\end{block}

\begin{exampleblock}{Flexibilidad}
Cada panel puede configurarse según las preferencias del expositor.
\end{exampleblock}

\column{0.38\textwidth}
\dashboardcard{Consejo}{Mantén una distribución consistente entre presentaciones.}{Reduce errores y mejora la fluidez al exponer.}
\end{columns}

\note{
No es obligatorio este layout, pero sí es el más natural.
}
\end{frame}

% =================================================
% Etiquetas
% =================================================
\begin{frame}{Etiquetas de panel}
\begin{columns}[T,totalwidth=\textwidth]
\column{0.56\textwidth}
\imgplaceholder{presenter-labels.png}{Paneles con etiquetas visibles para identificar rápidamente su contenido.}

\column{0.40\textwidth}
\begin{block}{Utilidad}
\begin{itemize}
  \item identificar cada panel;
  \item validar configuración;
  \item evitar errores antes de exponer.
\end{itemize}
\end{block}

\begin{alertblock}{Recomendación}
Activa las etiquetas al inicio y desactívalas antes de presentar.
\end{alertblock}
\end{columns}

\note{
Muy útil durante pruebas. No tanto durante exposición real.
}
\end{frame}

% =================================================
% Abrir PDF (visual)
% =================================================
\begin{frame}{Abrir un PDF}
\begin{columns}[T,totalwidth=\textwidth]
\column{0.56\textwidth}
\imgplaceholder{open-pdf.png}{Captura de la aplicación con un PDF recién abierto.}

\column{0.40\textwidth}
\begin{block}{Qué ocurre al abrir}
\begin{itemize}
  \item se carga el documento;
  \item se asignan paneles automáticamente;
  \item se detectan notes si existen;
  \item se habilita navegación.
\end{itemize}
\end{block}

\begin{alertblock}{Importante}
Verifica si el PDF fue interpretado correctamente (con o sin notes).
\end{alertblock}
\end{columns}

\note{
Este slide conecta la teoría con la acción.
}
\end{frame}

% =================================================
% Navegación conceptual
% =================================================
\begin{frame}{Navegación básica}
\begin{columns}[T,totalwidth=\textwidth]
\column{0.58\textwidth}
\begin{itemize}
  \item avanzar y retroceder entre páginas;
  \item revisar la siguiente filmina;
  \item saltar al inicio o final;
  \item mantener ritmo de exposición.
\end{itemize}

\begin{block}{Importancia}
La navegación debe ser \bluebf{rápida, natural y sin distracciones}.
\end{block}

\column{0.38\textwidth}
\dashboardcard{Buena práctica}{Ensaya navegación con el PDF real.}{No improvises el ritmo durante la presentación.}
\end{columns}

\note{
Aquí todavía no metemos atajos. Eso va en la siguiente sección.
}
\end{frame}

% =================================================
% Panel context menu
% =================================================
\begin{frame}{Configuración de paneles}
\begin{columns}[T,totalwidth=\textwidth]
\column{0.56\textwidth}
\imgplaceholder{panel-context-menu.png}{Menú contextual para cambiar el contenido de cada panel.}

\column{0.40\textwidth}
\begin{block}{Opciones típicas}
\begin{itemize}
  \item filmina actual;
  \item siguiente;
  \item notas actuales;
  \item notas siguientes;
  \item vistas auxiliares.
\end{itemize}
\end{block}
\end{columns}

\note{
Aquí empieza el poder real del sistema: personalización.
}
\end{frame}

% =================================================
% Layout persistente
% =================================================
\begin{frame}{Persistencia del layout}
\begin{columns}[T,totalwidth=\textwidth]
\column{0.58\textwidth}
\begin{block}{Archivo de configuración}
El visor guarda automáticamente la distribución de paneles en un archivo:

\vspace{0.3cm}
\codebox{layout.json}
\end{block}

\begin{itemize}
  \item almacena la configuración de paneles;
  \item guarda posiciones y contenidos;
  \item permite restaurar el layout al reiniciar.
\end{itemize}

\column{0.38\textwidth}
\dashboardcard{Ventaja}{No necesitas reconfigurar cada vez.}{El entorno de trabajo se mantiene consistente entre sesiones.}
\end{columns}

\note{
Este archivo se genera automáticamente después de modificar el layout.
}
\end{frame}

\begin{frame}{Ubicación de layout.json}
\begin{block}{Ruta típica}
El archivo se crea en el directorio de trabajo del usuario o en la carpeta del proyecto:

\vspace{0.3cm}
\codebox{./layout.json}
\end{block}

\begin{itemize}
  \item se genera automáticamente;
  \item se sobrescribe al cambiar el layout;
  \item puede eliminarse para reiniciar la configuración.
\end{itemize}

\begin{alertblock}{Consejo}
Si algo se rompe visualmente, eliminar este archivo puede restaurar el comportamiento por defecto.
\end{alertblock}

\note{
Muy útil para debugging.
}
\end{frame}

% =================================================
% Panel activo
% =================================================
\begin{frame}{Selección de panel}
\begin{block}{Panel activo}
El sistema permite trabajar con un \bluebf{panel seleccionado} para aplicar acciones específicas.
\end{block}

\begin{itemize}
  \item cambiar contenido de un panel;
  \item ajustar layout;
  \item trabajar con teclado de forma precisa.
\end{itemize}

\begin{exampleblock}{Ventaja}
Permite una configuración avanzada sin depender del mouse.
\end{exampleblock}

\note{
Esto conecta directo con la sección de atajos.
}
\end{frame}

% =================================================
% Proyector
% =================================================
\begin{frame}{Control del proyector}
\begin{block}{Concepto clave}
El presentador y la audiencia ven cosas diferentes:
\end{block}

\begin{columns}[T,totalwidth=\textwidth]
\column{0.48\textwidth}
\imgplaceholder{open-pdf.png}{Configuración con dos pantallas.}

\column{0.48\textwidth}
\imgplaceholder{projector-view.png}{Salida del proyector (lo que ve el público).}
\end{columns}

\begin{alertblock}{Importante}
El público nunca debe ver notas ni paneles auxiliares.
\end{alertblock}

\note{
Este slide vende el valor principal del proyecto.
}
\end{frame}

\section{Teclas y herramientas}
\sectioncard{3. Teclas y herramientas}{Control total durante la presentación}

% =================================================
% Navegación
% =================================================
\begin{frame}{Navegación con teclado}
\shortcutline{Avanzar}{\kbd{Right}\ \kbd{Space}\ \kbd{Down}\ \kbd{PageDown}}
\shortcutline{Retroceder}{\kbd{Left}\ \kbd{Backspace}\ \kbd{Up}\ \kbd{PageUp}}
\shortcutline{Inicio y final}{\kbd{Home}\ \kbd{End}}

\note{
Este es el mínimo indispensable para controlar una presentación sin mouse.
}
\end{frame}

% =================================================
% Tipos de panel
% =================================================
\begin{frame}{Tipos de panel}
\begin{columns}[T,totalwidth=\textwidth]
\column{0.56\textwidth}
\imgplaceholder{presenter-labels.png}{Ejemplo de los distintos tipos de panel disponibles.}

\column{0.40\textwidth}
\begin{block}{Paneles disponibles}
\begin{itemize}
  \item filmina actual;
  \item siguiente filmina;
  \item notas actuales;
  \item notas siguientes.
\end{itemize}
\end{block}

\begin{exampleblock}{Idea clave}
Cada panel puede mostrar contenido diferente al mismo tiempo.
\end{exampleblock}
\end{columns}

\note{
Aquí introduces mentalmente el concepto de tipos de panel.
}
\end{frame}

% =================================================
% Paneles duplicados
% =================================================
\begin{frame}{Uso de paneles duplicados}
\begin{block}{Comportamiento}
El visor permite asignar el mismo tipo de contenido a múltiples paneles.
\end{block}

\begin{itemize}
  \item ver varias copias de la filmina actual;
  \item mostrar notes en más de un panel;
  \item crear layouts personalizados;
  \item adaptar el entorno al estilo del expositor.
\end{itemize}

\begin{exampleblock}{Ejemplo}
Dos paneles pueden mostrar la filmina actual mientras otros muestran notes.
\end{exampleblock}

\note{
Esto no es común en todos los visores, es una ventaja fuerte.
}
\end{frame}

% =================================================
% Selección de panel
% =================================================
\begin{frame}{Selección de panel}
\shortcutline{Seleccionar panel}{\kbd{Ctrl/Cmd+1}\ \kbd{Ctrl/Cmd+2}\ \kbd{Ctrl/Cmd+3}\ \kbd{Ctrl/Cmd+4}}

\begin{block}{Panel activo}
Permite trabajar sobre un panel específico sin usar el mouse.
\end{block}

\begin{itemize}
  \item seleccionar rápidamente un panel;
  \item cambiar su contenido;
  \item configurar layout con precisión.
\end{itemize}

\note{
Esto conecta con flujo avanzado sin mouse.
}
\end{frame}

% =================================================
% Asignación de contenido
% =================================================
\begin{frame}{Asignar contenido a un panel}
\shortcutline{Asignar contenido}{\kbd{Alt/Option+1}\ \kbd{Alt/Option+2}\ \kbd{Alt/Option+3}\ \kbd{Alt/Option+4}}

\begin{block}{Qué hace}
Permite cambiar el tipo de contenido del panel seleccionado.
\end{block}

\begin{itemize}
  \item cambiar entre filmina, siguiente y notes;
  \item configurar layout rápidamente;
  \item evitar uso de menú contextual.
\end{itemize}

\begin{exampleblock}{Flujo típico}
Seleccionar panel $\rightarrow$ asignar contenido.
\end{exampleblock}

\note{
Este es el punto donde el sistema se vuelve realmente potente.
}
\end{frame}

% =================================================
% Personalización
% =================================================
\begin{frame}{Personalización del layout}
\shortcutline{Modo customize}{\kbd{Shift+C}}
\shortcutline{Mostrar etiquetas}{\kbd{L}}

\begin{block}{Qué permite}
\begin{itemize}
  \item reorganizar paneles;
  \item visualizar etiquetas;
  \item ajustar el entorno de trabajo.
\end{itemize}
\end{block}

\begin{alertblock}{Recomendación}
Configura el layout antes de la presentación, no durante.
\end{alertblock}

\note{
Esto conecta con layout.json.
}
\end{frame}

% =================================================
% Modos principales
% =================================================
\begin{frame}{Modos principales}
\shortcutline{Cambiar de modo}{\kbd{1} normal \hspace{0.4cm} \kbd{2} pointer \hspace{0.4cm} \kbd{3} pen \hspace{0.4cm} \kbd{4} eraser \hspace{0.4cm} \kbd{5} spotlight}

\begin{columns}[T,totalwidth=\textwidth]
\column{0.56\textwidth}
\begin{block}{Qué hace cada modo}
\begin{itemize}
  \item \textbf{Normal}: navegación general.
  \item \textbf{Pointer}: señalar elementos.
  \item \textbf{Pen}: dibujar.
  \item \textbf{Eraser}: borrar trazos.
  \item \textbf{Spotlight}: enfocar zona.
\end{itemize}
\end{block}

\column{0.40\textwidth}
\begin{exampleblock}{Idea clave}
Estos modos definen la interacción en tiempo real.
\end{exampleblock}
\end{columns}

\note{
Memorizar 1-5 es clave.
}
\end{frame}

% =================================================
% Pointer
% =================================================
\begin{frame}{Uso del puntero}
\begin{columns}[T,totalwidth=\textwidth]
\column{0.56\textwidth}
\imgplaceholder{pointer.png}{Uso del puntero sobre la filmina.}

\column{0.40\textwidth}
\begin{block}{Para qué sirve}
\begin{itemize}
  \item señalar elementos;
  \item guiar atención;
  \item enfatizar.
\end{itemize}
\end{block}

\begin{alertblock}{Evitar}
Movimiento constante sin intención.
\end{alertblock}
\end{columns}
\end{frame}

% =================================================
% Pen
% =================================================
\begin{frame}{Dibujo sobre la presentación}
\begin{columns}[T,totalwidth=\textwidth]
\column{0.56\textwidth}
\imgplaceholder{drawing-tools.png}{Dibujo activo sobre la filmina.}

\column{0.40\textwidth}
\begin{block}{Aplicaciones}
\begin{itemize}
  \item subrayar;
  \item explicar;
  \item resolver en vivo.
\end{itemize}
\end{block}

\begin{exampleblock}{Ajuste de tamaño}
Con \kbd{+} y \kbd{-} se puede aumentar o disminuir el grosor del trazo del \textbf{pen}.
\end{exampleblock}
\end{columns}

\shortcutline{Control}{\kbd{3} pen \hspace{0.4cm} \kbd{+/-} tamaño \hspace{0.4cm} \kbd{C} limpiar \hspace{0.4cm} \kbd{D} mostrar/ocultar}
\end{frame}

% =================================================
% Eraser
% =================================================
\begin{frame}{Uso del borrador}
\begin{columns}[T,totalwidth=\textwidth]
\column{0.56\textwidth}
\imgplaceholder{eraser.png}{Uso del borrador sobre trazos.}

\column{0.40\textwidth}
\begin{block}{Qué hace}
\begin{itemize}
  \item elimina trazos;
  \item permite correcciones;
  \item limpia zonas específicas.
\end{itemize}
\end{block}

\begin{alertblock}{Diferencia}
\textbf{Eraser} borra parcialmente.  
\textbf{\kbd{C}} limpia todo.
\end{alertblock}

\begin{exampleblock}{Ajuste de tamaño}
Con \kbd{+} y \kbd{-} también puede ajustarse el tamaño del \textbf{eraser}.
\end{exampleblock}
\end{columns}

\shortcutline{Modo}{\kbd{4} eraser \hspace{0.6cm} \kbd{+/-} tamaño}
\end{frame}

% =================================================
% Spotlight
% =================================================
\begin{frame}{Spotlight}
\begin{block}{Concepto}
Enfoca una zona y oscurece el resto.
\end{block}

\begin{itemize}
  \item ideal para explicar;
  \item guía la atención;
  \item mejora claridad visual.
\end{itemize}

\begin{alertblock}{Uso}
Aplicar solo cuando sea necesario.
\end{alertblock}
\end{frame}

% =================================================
% Zoom
% =================================================
\begin{frame}{Selección y zoom}
\begin{columns}[T,totalwidth=\textwidth]
\column{0.56\textwidth}
\imgplaceholder{zoom-selection.png}{Selección para zoom.}

\column{0.40\textwidth}
\begin{block}{Flujo}
\begin{itemize}
  \item seleccionar;
  \item aplicar zoom;
  \item salir con \kbd{Esc}.
\end{itemize}
\end{block}
\end{columns}

\shortcutline{Zoom}{Seleccionar \hspace{0.4cm} \kbd{Z} aplicar \hspace{0.4cm} \kbd{Esc}}
\end{frame}

% =================================================
% Cronómetro
% =================================================
\begin{frame}{Cronómetro}
\shortcutline{Control}{\kbd{T} iniciar/pausar \hspace{0.8cm} \kbd{Shift+T} reiniciar}

\begin{itemize}
  \item controlar tiempo;
  \item ajustar ritmo;
  \item evitar excederse.
\end{itemize}

\begin{exampleblock}{Consejo}
Ensayar con cronómetro activo.
\end{exampleblock}
\end{frame}

% =================================================
% Proyector
% =================================================
\begin{frame}{Control del proyector}
\shortcutline{Black}{\kbd{B}}
\shortcutline{Freeze}{\kbd{F}}

\begin{itemize}
  \item ocultar pantalla;
  \item congelar imagen;
  \item controlar atención.
\end{itemize}
\end{frame}

% =================================================
% Interfaz general
% =================================================
\begin{frame}{Interfaz y control}
\shortcutline{Ayuda}{\kbd{H}}
\shortcutline{Fullscreen}{\kbd{W}}
\shortcutline{Salir}{\kbd{Esc}}
\shortcutline{Pnotes}{\kbd{P}}
\end{frame}

% =================================================
% Resumen
% =================================================
\begin{frame}{Resumen de atajos clave}
\begin{block}{Memoriza esto}
\centering

\kbd{← →} \hspace{0.35cm}
\kbd{1-5} \hspace{0.35cm}
\kbd{+/-} \hspace{0.35cm}
\kbd{T} \hspace{0.35cm}
\kbd{B} \hspace{0.35cm}
\kbd{F} \hspace{0.35cm}
\kbd{Esc}

\end{block}

\begin{exampleblock}{Idea final}
Puedes presentar completamente sin usar el mouse.
\end{exampleblock}
\end{frame}

\section{Preparación del PDF}
\sectioncard{4. Preparación del PDF}{Configuración correcta desde LaTeX}

% =================================================
% Introducción
% =================================================
\begin{frame}{Preparación desde LaTeX}
\begin{block}{Idea principal}
Para aprovechar completamente Presenter Viewer, el PDF debe prepararse correctamente desde LaTeX.
\end{block}

\begin{itemize}
  \item uso de Beamer;
  \item separación entre contenido y notas;
  \item compatibilidad con notes y pnotes;
  \item exportación adecuada del PDF.
\end{itemize}

\begin{exampleblock}{Clave}
El visor funciona mejor cuando el PDF fue diseñado pensando en el presentador.
\end{exampleblock}

\note{
Aquí introduces que no es solo abrir cualquier PDF.
}
\end{frame}

% =================================================
% Cabecera mínima
% =================================================
\begin{frame}[fragile]{Configuración mínima recomendada}
\begin{block}{Cabecera esencial}
\begin{lstlisting}
\documentclass[aspectratio=169]{beamer}
\usepackage{pdfcomment}
\setbeameroption{show notes on second screen=left}
\end{lstlisting}
\end{block}

\begin{itemize}
  \item activa exportación de notes;
  \item genera layout compatible con visor;
  \item separa slide y notas en el PDF.
\end{itemize}

\begin{alertblock}{Importante}
Sin esta configuración, las notes no aparecerán como panel separado.
\end{alertblock}

\note{
Esto es obligatorio para aprovechar notes.
}
\end{frame}

% =================================================
% Explicación visual
% =================================================
\begin{frame}{¿Qué hace esta configuración?}
\begin{block}{Resultado en el PDF}
\begin{itemize}
  \item lado izquierdo: notes;
  \item lado derecho: filmina;
  \item layout ancho tipo presentador.
\end{itemize}
\end{block}

\begin{exampleblock}{Interpretación}
El visor puede separar automáticamente ambas partes.
\end{exampleblock}

\note{
Esto conecta directamente con cómo el visor interpreta el PDF.
}
\end{frame}

% =================================================
% Uso de note
% =================================================
\begin{frame}{Uso de \texttt{\textbackslash note\{\}}}
\begin{block}{Ejemplo}
\codebox{\textbackslash begin\{frame\}\{Mi filmina\}\\
Contenido visible para el publico.\\
\textbackslash note\{Esta nota solo la ve el presentador.\}}
\end{block}

\begin{itemize}
  \item no se muestra al público;
  \item aparece en panel de notes;
  \item sirve como guía del expositor.
\end{itemize}

\begin{exampleblock}{Ventaja}
Permite tener apoyo sin saturar la diapositiva.
\end{exampleblock}

\note{
Este es el mecanismo principal.
}
\end{frame}

% =================================================
% Ejemplo conceptual
% =================================================
\begin{frame}{Ejemplo con notes}
\begin{columns}[T,totalwidth=\textwidth]
\column{0.56\textwidth}
\begin{block}{Filmina}
Pregunta: \bluebf{¿Por qué usar un visor de presentador?}
\end{block}

\begin{itemize}
  \item control del tiempo;
  \item notas privadas;
  \item herramientas visuales;
  \item mejor flujo de exposición.
\end{itemize}

\column{0.40\textwidth}
\begin{exampleblock}{Qué sucede}
El público ve esta filmina, pero el expositor ve también las notes.
\end{exampleblock}
\end{columns}

\note{
Respuesta sugerida: mejora control, claridad y ritmo.
}
\end{frame}

% =================================================
% pnote macro
% =================================================
\begin{frame}[fragile]{Agregar \texttt{\textbackslash pnote}}
\begin{block}{Definición}
\begin{lstlisting}
\newcommand{\pnote}[1]{%
  \pdfcomment[
    icon=Note,
    open=false,
    subject={Presenter Note},
    color={1 1 0},
    opacity=0
  ]{#1}%
}
\end{lstlisting}
\end{block}

\begin{alertblock}{Importante}
Debe agregarse en el preámbulo del archivo.
\end{alertblock}

\note{
Esto habilita notas embebidas en el PDF.
}
\end{frame}

% =================================================
% Uso de pnote
% =================================================
\begin{frame}{Uso de \texttt{\textbackslash pnote\{\}}}
\begin{block}{Ejemplo}
\codebox{Explicación importante.\\
\textbackslash pnote\{Recordar dar ejemplo practico.\}}
\end{block}

\begin{itemize}
  \item recordatorios breves;
  \item notas tácticas;
  \item apoyo rápido durante exposición.
\end{itemize}

\begin{exampleblock}{Idea}
Complementa a \texttt{\textbackslash note}, no lo reemplaza.
\end{exampleblock}

\note{
pnote es más ligero y puntual.
}
\end{frame}

% =================================================
% Comparación
% =================================================
\begin{frame}{Notes vs Pnotes}
\begin{table}
\centering
\small
\begin{tabular}{p{3cm} p{4.5cm} p{4.5cm}}
\toprule
Método & Ventajas & Uso \\
\midrule
\texttt{\textbackslash note} &
Integración completa con Beamer &
Notas amplias y estructuradas \\
\addlinespace
\texttt{\textbackslash pnote} &
Ligero y flexible &
Recordatorios breves \\
\bottomrule
\end{tabular}
\end{table}

\begin{exampleblock}{Recomendación}
Usar ambos en conjunto.
\end{exampleblock}

\note{
Esta es la decisión de diseño recomendada.
}
\end{frame}

% =================================================
% PDF sin notes
% =================================================
\begin{frame}{PDFs sin panel de notas}
\begin{alertblock}{Comportamiento esperado}
Si el PDF no tiene notes, el visor debe mostrar la página completa.
\end{alertblock}

\begin{itemize}
  \item no recortar contenido;
  \item mostrar PDF normal;
  \item seguir permitiendo pnotes.
\end{itemize}

\begin{exampleblock}{Importante}
No todos los PDFs requieren notes.
\end{exampleblock}

\note{
Esto es un edge case importante.
}
\end{frame}

% =================================================
% Buenas prácticas
% =================================================
\begin{frame}{Buenas prácticas}
\begin{itemize}
  \item una idea por diapositiva;
  \item evitar saturar texto;
  \item usar notes para apoyo;
  \item priorizar claridad visual;
  \item ensayar con el visor.
\end{itemize}

\begin{alertblock}{Clave}
Un buen PDF es tan importante como el visor.
\end{alertblock}

\note{
Cerrar con enfoque pedagógico.
}
\end{frame}

\section{Buenas prácticas}
\sectioncard{5. Buenas prácticas}{Diseño, preparación y exposición}

% =================================================
% Diseño
% =================================================
\begin{frame}{Buenas prácticas de diseño}
\begin{itemize}
  \item una idea principal por diapositiva;
  \item texto mínimo y bien espaciado;
  \item usar contraste adecuado;
  \item priorizar lo visual sobre lo textual;
  \item evitar saturar la slide.
\end{itemize}

\begin{exampleblock}{Regla simple}
Si una diapositiva tiene demasiado texto, probablemente debe dividirse.
\end{exampleblock}

\note{
Esto conecta con principios básicos de presentación.
}
\end{frame}

% =================================================
% Uso del visor
% =================================================
\begin{frame}{Buenas prácticas con el visor}
\begin{itemize}
  \item usar pointer solo cuando sea necesario;
  \item usar spotlight para guiar la atención;
  \item evitar cambiar constantemente de modo;
  \item limpiar dibujos cuando ya no se necesiten;
  \item mantener consistencia en el layout.
\end{itemize}

\begin{alertblock}{Clave}
El visor debe apoyar la exposición, no distraer.
\end{alertblock}

\note{
Aquí se aterriza el uso correcto de las herramientas.
}
\end{frame}

% =================================================
% Preparación
% =================================================
\begin{frame}{Preparación antes de exponer}
\begin{itemize}
  \item ensayar con el PDF final;
  \item validar notes y pnotes;
  \item configurar paneles;
  \item probar herramientas;
  \item verificar el proyector.
\end{itemize}

\begin{exampleblock}{Consejo}
No pruebes cosas nuevas en vivo.
\end{exampleblock}

\note{
Esto es experiencia real de exposición.
}
\end{frame}

% =================================================
% Checklist
% =================================================
\begin{frame}{Checklist antes de iniciar}
\begin{block}{Verificación rápida}
\begin{itemize}
  \item PDF listo;
  \item Presenter Viewer abierto;
  \item paneles configurados;
  \item proyector funcionando;
  \item notes visibles correctamente;
  \item cronómetro listo.
\end{itemize}
\end{block}

\begin{alertblock}{Importante}
Haz esta revisión antes de estar frente al público.
\end{alertblock}


\end{frame}

\section{Flujo recomendado}
\sectioncard{6. Flujo recomendado}{Desde LaTeX hasta la exposición}

% =================================================
% Flujo completo
% =================================================
\begin{frame}{Flujo completo recomendado}
\begin{block}{Pipeline de trabajo}
\centering
\kbd{LaTeX} \hspace{0.35cm} $\rightarrow$ \hspace{0.35cm}
\kbd{PDF} \hspace{0.35cm} $\rightarrow$ \hspace{0.35cm}
\kbd{Presenter Viewer} \hspace{0.35cm} $\rightarrow$ \hspace{0.35cm}
\kbd{Configurar} \hspace{0.35cm} $\rightarrow$ \hspace{0.35cm}
\kbd{Presentar}
\end{block}

\vspace{0.3cm}

\begin{itemize}
  \item crear la presentación en Beamer;
  \item exportar a PDF;
  \item abrir en Presenter Viewer;
  \item configurar paneles y herramientas;
  \item realizar la exposición.
\end{itemize}

\note{
Este slide resume todo el flujo en una sola vista clara.
}
\end{frame}

% =================================================
% Flujo detallado
% =================================================
\begin{frame}{Flujo paso a paso}
\begin{enumerate}
    \item crear la presentación en LaTeX (Beamer);
    \item agregar \texttt{\textbackslash note\{\}} para notas extensas;
    \item agregar la macro \texttt{\textbackslash pnote\{\}} al preámbulo;
    \item insertar pnotes donde se necesiten recordatorios rápidos;
    \item compilar el PDF final;
    \item abrir el archivo en Presenter Viewer;
    \item configurar layout y paneles;
    \item probar herramientas (pointer, pen, zoom, etc.);
    \item ensayar la presentación completa.
\end{enumerate}

\begin{alertblock}{Consejo práctico}
Haz al menos un ensayo completo antes de la presentación real.
\end{alertblock}

\note{
Este es el flujo real que deberían seguir los usuarios.
}
\end{frame}

% =================================================
% Errores comunes
% =================================================
\begin{frame}{Errores comunes}
\begin{itemize}
  \item no probar el PDF antes de exponer;
  \item depender solo de memoria (sin notes);
  \item saturar slides con demasiado texto;
  \item abusar de herramientas visuales;
  \item no verificar el proyector o segunda pantalla;
  \item improvisar configuración en vivo.
\end{itemize}

\begin{alertblock}{Importante}
La mayoría de errores no son técnicos, son de preparación.
\end{alertblock}

\note{
Este slide agrega mucho valor práctico.
}
\end{frame}

% =================================================
% Cierre
% =================================================
\begin{frame}{Cierre}
\begin{block}{Idea principal}
Presenter Viewer es una herramienta que mejora significativamente la calidad de una exposición cuando se usa correctamente.
\end{block}

\begin{itemize}
  \item mejora el control del expositor;
  \item permite separar contenido público y privado;
  \item facilita el uso de herramientas visuales;
  \item ayuda a mantener ritmo y claridad.
\end{itemize}

\begin{exampleblock}{Clave}
La herramienta no reemplaza una buena presentación, pero la potencia.
\end{exampleblock}

\note{
Cierre conceptual fuerte.
}
\end{frame}

% =================================================
% Referencia
% =================================================
\begin{frame}{Referencia de inspiración}
La organización funcional de esta herramienta toma inspiración de soluciones como \texttt{pdfpc}, especialmente en:

\begin{itemize}
    \item vista del presentador;
    \item manejo de notas;
    \item pointer y spotlight;
    \item herramientas de dibujo;
    \item cronómetro;
    \item control del proyector.
\end{itemize}

\soft{La intención es ofrecer una versión moderna, extensible y adaptada a flujos actuales.}

\note{
Esto posiciona el proyecto sin verse como copia.
}
\end{frame}

\end{document}
